\section{Consistency Models}
\label{sec:consistency-models}

When users access the same data, what guarantees can the application provide about what they see?
A \emph{consistency model} specifies permitted observations: when writes become visible, how operations are ordered, and whether replicas agree after receiving the same updates.
\parencite[Section~10.3]{understanding-distributed-systems}

Neither architecture guarantees a particular consistency model on its own.
A server-centred application can require coordination before completing an operation, whereas a local-first application must also allow supported operations to complete while disconnected.

\subsection{Strong consistency and current shared state}

Following \textcite[85]{understanding-distributed-systems}, \emph{strong consistency} means that operations behave as if there were a single copy of the data.
Each operation appears to take effect at one instant between its invocation and completion.
Once a write completes, its effects must be reflected in subsequent reads.\footnote{Vitillo identifies this guarantee as \emph{linearizability} on page~86.
	This chapter uses \emph{strong consistency} in that sense.}

In the server-centred task board, Alice changes a title from \enquote{Draft} to \enquote{Review}.
If Bob reads it after her save completes, with no intervening writes, he must see \enquote{Review} (Figure~\ref{fig:strong-consistency-timeline}).
The bars span each operation's invocation and completion.

\begin{illustration}[htbp]
	\centering
	\caption{A read following a completed write under strong consistency, assuming no intervening writes}
	\label{fig:strong-consistency-timeline}
	\begingroup
	\setstretch{1}
	\begin{tikzpicture}[thesis]
	\draw[thesis/local, ->] (1.7,0) -- (13.2,0) node[above left] {Time};
	\node[anchor=east] at (1.3,-1.2) {Alice};
	\node[anchor=east] at (1.3,-2.4) {Bob};
	\draw[thesis/lifeline] (1.7,-1.2) -- (13.2,-1.2);
	\draw[thesis/lifeline] (1.7,-2.4) -- (13.2,-2.4);

	\draw[thesis/box] (2.2,-1.55) rectangle (6.6,-0.85);
	\node at (4.4,-1.2) {Write title = \enquote{Review}};
	\draw[thesis/divider] (6.6,-0.45) -- (6.6,-2.95);
	\fill (6.6,-1.2) circle (1.8pt);
	\node[thesis/label, above] at (6.6,-0.45) {Write completed};

	\draw[thesis/box] (8.5,-2.75) rectangle (11.3,-2.05);
	\node at (9.9,-2.4) {Read title};
	\fill (8.5,-2.4) circle (1.8pt);
	\node[thesis/label, above] at (8.5,-1.95) {Read begins later};
	\draw[thesis/local, ->] (11.3,-2.75) -- (11.3,-3.1);
	\node[below] at (11.3,-3.1) {Must return \enquote{Review}};
\end{tikzpicture}

	\endgroup
	\figuresource{Author's illustration based on \textcite[85]{understanding-distributed-systems} and \textcite[324--328]{designing-data-intensive-applications}.}
\end{illustration}
\FloatBarrier

The guarantee must cover the entire read path, including caches and database replicas.
\parencite[Section~10.3.1]{understanding-distributed-systems}

\newpage
\subsection{CAP theorem and offline availability}

In the local-first version, Alice saves \enquote{Review} locally on her laptop while Bob's phone is disconnected.
Bob's phone still holds \enquote{Draft}.
Bob's phone can answer the read using its local copy, but that copy does not include Alice's completed write.
Obtaining the current shared value requires communication, so the read cannot complete independently offline.
This is the trade-off illustrated through a partitioned key-value store by \textcite[Chapter~6]{system-design-interview}.
The \emph{CAP theorem} makes the limitation precise, using the following terms:
\begin{itemize}
	\item \emph{Consistency (C)} means strong consistency as defined above: operations behave as if performed on one copy, respecting real-time order.
	\item \emph{Availability (A)} requires every request received by a non-failing node to eventually return a valid operation result.
	\item \emph{Partition tolerance (P)} requires the guarantees to hold even when communication between parts of the system fails.
\end{itemize}
The CAP theorem states that during a network partition, a system cannot guarantee both strong consistency and availability.
Preserving strong consistency may require delaying or rejecting operations.
Keeping operations available on both sides of the partition instead means accepting that reads may return outdated values.
Weaker guarantees, such as eventual consistency, remain possible.
\parencites[336--338]{designing-data-intensive-applications}[Section~10.3.4]{understanding-distributed-systems}

Local-first applications favour availability during disconnection, whereas server-centred applications can require coordination to preserve strong consistency.
Neither architecture automatically determines the system's CAP guarantees.
Even after communication resumes, synchronization takes time; \emph{eventual} does not specify a maximum delay.
\parencite[Section~10.3]{understanding-distributed-systems}

\subsection{Eventual and strong eventual consistency}

Replicas may temporarily hold different data and still converge as updates are exchanged.
Eventual consistency describes this guarantee in both server-centred and local-first applications.

\textcite[Section~2.2, Definition~2]{conflict-free-replicated-data-types} define \emph{eventual consistency} through three properties:
\begin{itemize}
	\item \emph{Eventual delivery}: an update delivered at one correct replica eventually reaches all correct replicas.
	\item \emph{Convergence}: correct replicas that have delivered the same updates eventually reach equivalent state.
	\item \emph{Termination}: all method executions terminate.
\end{itemize}
A correct replica is one that has not crashed, and equivalent states return the same results to queries.

\pagebreak
\emph{Strong eventual consistency} (SEC) strengthens the convergence requirement as follows \parencite[Section~2.2, Definition~3]{conflict-free-replicated-data-types}:
\begin{quote}
	Correct replicas that have delivered the same updates have equivalent state.
\end{quote}

Replicas have equivalent states as soon as they have delivered the same updates, without further reconciliation.
Here, the word \emph{strong} refers to convergence; SEC does not provide the freshness guarantee of the strong consistency model.
Replicas can still hold different states while each has updates the other has not received.

Figure~\ref{fig:eventual-consistency-convergence} compares the guarantees after both replicas have delivered the same updates, assuming no further updates arrive.
Temporary disagreement is permitted under EC, but is not required.

\begin{illustration}[htbp]
	\centering
	\caption{Convergence after two correct replicas have delivered the same updates}
	\label{fig:eventual-consistency-convergence}
	\begingroup
	\setstretch{1}
	\begin{tikzpicture}[
	thesis,
	state/.style={thesis/box, text width=3.4cm, minimum height=0.8cm, inner sep=4pt}
]
	\node[thesis/label] at (6.5,0) {Same updates\\delivered at both replicas};
	\node[thesis/label] at (11.7,0) {Later\\no new updates};
	\node[anchor=east, align=right] at (4.1,-1.1) {Eventual\\consistency};
	\node[anchor=east, align=right] at (4.1,-2.25) {Strong eventual\\consistency};

	\node[state] (ec-start) at (6.5,-1.1) {States may differ};
	\node[state] (ec-end) at (11.7,-1.1) {Equivalent states};
	\draw[thesis/local, ->] (ec-start.east) -- (ec-end.west);
	\node[state] (sec-start) at (6.5,-2.25) {Equivalent states};
	\node[state] (sec-end) at (11.7,-2.25) {Equivalent states};
	\draw[thesis/local, ->] (sec-start.east) -- (sec-end.west);

	\node[align=center] at (8,-3.15) {Neither model guarantees that all recent updates have arrived.};
\end{tikzpicture}

	\endgroup
	\figuresource{Author's illustration based on \textcite[Section~2.2, Definitions~2--3]{conflict-free-replicated-data-types}.}
\end{illustration}
\FloatBarrier

Local-first applications allow disconnected devices to accept changes independently.
Synchronization must reconcile these changes when the devices reconnect.
\parencite[Chapter~5]{designing-data-intensive-applications}

CRDTs provide sufficient conditions for SEC through their update and merge rules, assuming eventual delivery and termination.
This makes SEC particularly useful for local-first software: independently editable replicas can agree once they have received the same updates, without requiring server approval for each edit.
\parencite[Section~2.3, Theorem~1]{conflict-free-replicated-data-types}

Local-first architecture alone does not guarantee SEC.
Convergence also does not ensure that the merged result preserves every user's intention or every business rule.
\parencite[Section~11.2]{understanding-distributed-systems}
